% Appendix: Accessing the Server and Running the ISR Analysis Pipeline
% Include in your main document with:  \appendix  % Appendix: Accessing the Server and Running the ISR Analysis Pipeline
% Include in your main document with:  \appendix  % Appendix: Accessing the Server and Running the ISR Analysis Pipeline
% Include in your main document with:  \appendix  % Appendix: Accessing the Server and Running the ISR Analysis Pipeline
% Include in your main document with:  \appendix  \input{appendix_analysis_guide}

\section{Accessing the Analysis Server and Running the ISR Pipeline}
\label{app:analysis_guide}

This appendix describes how to obtain an account on the analysis server,
access a remote desktop session, retrieve the analysis code, set up the
Python environment, and run the Millstone Hill ISR plasma-parameter pipeline.

%-------------------------------------------------------------------
\subsection{Getting an Account on Revontuli}

The analysis is performed on \texttt{revontuli.uit.no}, a Linux server at
UiT -- The Arctic University of Norway.
Contact the group supervisor to have an account created.
You will receive a username and an initial password.
Change the password after first login:
\begin{verbatim}
ssh <username>@revontuli.uit.no
passwd
\end{verbatim}

%-------------------------------------------------------------------
\subsection{Remote Desktop Access}

A graphical desktop (XFCE4) is available via VNC.
Each user's VNC port is determined by their Unix user ID:
\begin{equation}
  \text{port} = \text{UID} + 50000
\end{equation}
Ask the system administrator for your UID, or read it yourself after
logging in with \texttt{id -u}.

\subsubsection*{Starting the desktop}

Log in over SSH and run the \texttt{remotedesktop} command (defined
system-wide in \texttt{/etc/profile.d/remotedesktop.sh}):
\begin{verbatim}
ssh <username>@revontuli.uit.no
remotedesktop
\end{verbatim}
This starts an Xtightvnc session with XFCE4 on your personal port.
To reconnect after a reboot, simply run \texttt{remotedesktop} again.

\subsubsection*{Connecting from your laptop}

Forward the VNC port over SSH so the connection is encrypted:
\begin{verbatim}
ssh -L 5901:localhost:<your-vnc-port> <username>@revontuli.uit.no
\end{verbatim}
Then open any VNC client (e.g.\ Remmina, TigerVNC, RealVNC) and connect
to \texttt{localhost:5901}.
The desktop appears without requiring a VNC password.

%-------------------------------------------------------------------
\subsection{Getting the Analysis Code}

The pipeline lives in a Git repository.
Clone it into your home directory on Revontuli:
\begin{verbatim}
cd ~
git clone https://github.com/jvierine/isr_analysis.git src/antistarlink
cd src/antistarlink
\end{verbatim}
To update to the latest version at any time:
\begin{verbatim}
git pull
\end{verbatim}

%-------------------------------------------------------------------
\subsection{Setting up the Python Environment}

A setup script installs all required Python packages into a virtual
environment at \texttt{\textasciitilde/venv/antistarlink}.
Run it once per machine:
\begin{verbatim}
cd ~/src/antistarlink
bash setup_env.sh
\end{verbatim}
The script requires \texttt{libfftw3-dev} to build \texttt{pyfftw}
(already installed on Revontuli).
It installs:
\begin{itemize}
  \item \texttt{pyfftw} -- fast FFT via FFTW3
  \item \texttt{digital\_rf} -- reading raw DigitalRF voltage data
  \item \texttt{mpi4py} -- MPI parallelisation
\end{itemize}
and inherits \texttt{h5py}, \texttt{scipy}, \texttt{numpy}, and
\texttt{matplotlib} from the system Python~3.14 installation.

%-------------------------------------------------------------------
\subsection{Raw Data Layout}

Each experiment directory contains the raw DigitalRF data and radar metadata:
\begin{verbatim}
/mnt/data/juha/millstone_hill/isr/<date>/
  usrp-rx0-r_<start>_<end>/
    rf_data/
      zenith-l/        # receive channel (1 MHz, complex64)
      misa-l/          # MISA receive channel (when available)
      tx-h/            # transmitted waveform copy
    metadata/
      id_metadata/     # pulse sweep IDs (coded=1-32, LP=300, scan=800)
      powermeter/      # transmit power vs. time
      antenna_control_metadata/   # pointing (azimuth, elevation)
\end{verbatim}

%-------------------------------------------------------------------
\subsection{Running the Pipeline}

The driver script \texttt{run\_analysis.py} reads a JSON configuration file
and executes the enabled pipeline steps in order.
Launch it with MPI to process intervals in parallel:
\begin{verbatim}
cd ~/src/antistarlink
mpirun -np 24 ~/venv/antistarlink/bin/python3 \
    run_analysis.py config/millstone_2023-09-05.json
\end{verbatim}
Replace \texttt{24} with the number of CPU cores available
(\texttt{nproc} reports the total on the current machine).

\subsubsection*{Configuration file}

An annotated example for the 2023-09-05 dataset is given below.
Create a copy and edit it for a new experiment.

\begin{verbatim}
{
  "experiment": "millstone_2023-09-05",
  "data_dir":   "/mnt/data/juha/millstone_hill/isr/2023-09-05/
                  usrp-rx0-r_20230905T214448_20230906T040054",
  "output_dir": "/mnt/data/juha/millstone_hill/results/2023-09-05",
  "max_time_s": 600,

  "steps": {
    "lpi": {
      "enabled":            true,
      "channel":            "zenith-l",
      "range_gate_us":      60,
      "avg_dur":            10,
      "pass_band_hz":       100000,
      "filter_len":         20,
      "max_range_delay_us": 7000,
      "save_acf_images":    true,
      "reanalyze":          false
    },
    "fit_lpi": {
      "enabled":   true,
      "channel":   "zenith-l",
      "max_dt":    300,
      "range_avg": [1, 3, 5],
      "reanalyze": false
    },
    "long_pulse": {
      "enabled": false,
      "channel": "zenith-l",
      "mode":    300
    },
    "fit_lp": {
      "enabled": false,
      "channel": "zenith-l",
      "avg_dur": 600,
      "ridx":    [35, 230]
    }
  }
}
\end{verbatim}

\subsubsection*{Key parameters to adjust}

\begin{description}
  \item[\texttt{channel}]
    Use \texttt{"zenith-l"} for the zenith antenna or \texttt{"misa-l"}
    for the steerable MISA antenna.
    Check which channels are present in \texttt{rf\_data/} before running.
  \item[\texttt{range\_gate\_us}]
    Range resolution in microseconds; must match the radar pulse-bit length.
    Typical values: 30, 60, or 120~\textmu s.
  \item[\texttt{max\_dt}]
    Integration time for the plasma-parameter fit.
    300~s is a good starting point; increase to 600~s for lower-SNR data.
  \item[\texttt{max\_time\_s}]
    Set to e.g.\ 600 for a quick 10-minute test run; remove the key
    (or set to \texttt{null}) to process the full experiment.
  \item[\texttt{reanalyze}]
    Set to \texttt{false} to skip intervals already on disk.
    Set to \texttt{true} to reprocess from scratch.
\end{description}

%-------------------------------------------------------------------
\subsection{Output Files}

All results are written to \texttt{output\_dir} (or inside the raw-data
directory if \texttt{output\_dir} is not set):

\begin{verbatim}
lpi_<range_gate_us>/zenith-l/
  lpi-<unix_timestamp>.h5    # ACF per range gate and lag (step 1)
  lpi-<unix_timestamp>.png   # diagnostic image

lpi_<range_gate_us>/zenith-l/
  pp-<unix_timestamp>.h5     # Te, Ti, vi, ne profiles (step 2)
  pp-<unix_timestamp>.png    # diagnostic plot

range_doppler_<mode>_outlier/zenith-l/
  il-<unix_timestamp>.h5     # range-Doppler spectra (step 3)
  pp-<unix_timestamp>.h5     # Te, Ti, vi, ne profiles (step 4)
\end{verbatim}

The plasma-parameter HDF5 files contain the following fields:
\begin{center}
\begin{tabular}{ll}
\hline
\textbf{Dataset} & \textbf{Content} \\
\hline
\texttt{Te}, \texttt{Ti}     & Electron / ion temperature (K) \\
\texttt{vi}                  & Line-of-sight ion drift velocity (m/s) \\
\texttt{ne}                  & Relative electron density proxy \\
\texttt{dTe/Ti}, \texttt{dTi}, \texttt{dvi}, \texttt{dne} & $1\sigma$ uncertainties \\
\texttt{rgs}                 & Range gate centres (km) \\
\texttt{t0}, \texttt{t1}     & Integration window (Unix seconds) \\
\texttt{az}, \texttt{el}     & Mean antenna pointing (degrees) \\
\texttt{T\_sys}              & System temperature (K) \\
\texttt{P\_tx}               & Mean transmit power (W) \\
\hline
\end{tabular}
\end{center}
% end of appendix


\section{Accessing the Analysis Server and Running the ISR Pipeline}
\label{app:analysis_guide}

This appendix describes how to obtain an account on the analysis server,
access a remote desktop session, retrieve the analysis code, set up the
Python environment, and run the Millstone Hill ISR plasma-parameter pipeline.

%-------------------------------------------------------------------
\subsection{Getting an Account on Revontuli}

The analysis is performed on \texttt{revontuli.uit.no}, a Linux server at
UiT -- The Arctic University of Norway.
Contact the group supervisor to have an account created.
You will receive a username and an initial password.
Change the password after first login:
\begin{verbatim}
ssh <username>@revontuli.uit.no
passwd
\end{verbatim}

%-------------------------------------------------------------------
\subsection{Remote Desktop Access}

A graphical desktop (XFCE4) is available via VNC.
Each user's VNC port is determined by their Unix user ID:
\begin{equation}
  \text{port} = \text{UID} + 50000
\end{equation}
Ask the system administrator for your UID, or read it yourself after
logging in with \texttt{id -u}.

\subsubsection*{Starting the desktop}

Log in over SSH and run the \texttt{remotedesktop} command (defined
system-wide in \texttt{/etc/profile.d/remotedesktop.sh}):
\begin{verbatim}
ssh <username>@revontuli.uit.no
remotedesktop
\end{verbatim}
This starts an Xtightvnc session with XFCE4 on your personal port.
To reconnect after a reboot, simply run \texttt{remotedesktop} again.

\subsubsection*{Connecting from your laptop}

Forward the VNC port over SSH so the connection is encrypted:
\begin{verbatim}
ssh -L 5901:localhost:<your-vnc-port> <username>@revontuli.uit.no
\end{verbatim}
Then open any VNC client (e.g.\ Remmina, TigerVNC, RealVNC) and connect
to \texttt{localhost:5901}.
The desktop appears without requiring a VNC password.

%-------------------------------------------------------------------
\subsection{Getting the Analysis Code}

The pipeline lives in a Git repository.
Clone it into your home directory on Revontuli:
\begin{verbatim}
cd ~
git clone https://github.com/jvierine/isr_analysis.git src/antistarlink
cd src/antistarlink
\end{verbatim}
To update to the latest version at any time:
\begin{verbatim}
git pull
\end{verbatim}

%-------------------------------------------------------------------
\subsection{Setting up the Python Environment}

A setup script installs all required Python packages into a virtual
environment at \texttt{\textasciitilde/venv/antistarlink}.
Run it once per machine:
\begin{verbatim}
cd ~/src/antistarlink
bash setup_env.sh
\end{verbatim}
The script requires \texttt{libfftw3-dev} to build \texttt{pyfftw}
(already installed on Revontuli).
It installs:
\begin{itemize}
  \item \texttt{pyfftw} -- fast FFT via FFTW3
  \item \texttt{digital\_rf} -- reading raw DigitalRF voltage data
  \item \texttt{mpi4py} -- MPI parallelisation
\end{itemize}
and inherits \texttt{h5py}, \texttt{scipy}, \texttt{numpy}, and
\texttt{matplotlib} from the system Python~3.14 installation.

%-------------------------------------------------------------------
\subsection{Raw Data Layout}

Each experiment directory contains the raw DigitalRF data and radar metadata:
\begin{verbatim}
/mnt/data/juha/millstone_hill/isr/<date>/
  usrp-rx0-r_<start>_<end>/
    rf_data/
      zenith-l/        # receive channel (1 MHz, complex64)
      misa-l/          # MISA receive channel (when available)
      tx-h/            # transmitted waveform copy
    metadata/
      id_metadata/     # pulse sweep IDs (coded=1-32, LP=300, scan=800)
      powermeter/      # transmit power vs. time
      antenna_control_metadata/   # pointing (azimuth, elevation)
\end{verbatim}

%-------------------------------------------------------------------
\subsection{Running the Pipeline}

The driver script \texttt{run\_analysis.py} reads a JSON configuration file
and executes the enabled pipeline steps in order.
Launch it with MPI to process intervals in parallel:
\begin{verbatim}
cd ~/src/antistarlink
mpirun -np 24 ~/venv/antistarlink/bin/python3 \
    run_analysis.py config/millstone_2023-09-05.json
\end{verbatim}
Replace \texttt{24} with the number of CPU cores available
(\texttt{nproc} reports the total on the current machine).

\subsubsection*{Configuration file}

An annotated example for the 2023-09-05 dataset is given below.
Create a copy and edit it for a new experiment.

\begin{verbatim}
{
  "experiment": "millstone_2023-09-05",
  "data_dir":   "/mnt/data/juha/millstone_hill/isr/2023-09-05/
                  usrp-rx0-r_20230905T214448_20230906T040054",
  "output_dir": "/mnt/data/juha/millstone_hill/results/2023-09-05",
  "max_time_s": 600,

  "steps": {
    "lpi": {
      "enabled":            true,
      "channel":            "zenith-l",
      "range_gate_us":      60,
      "avg_dur":            10,
      "pass_band_hz":       100000,
      "filter_len":         20,
      "max_range_delay_us": 7000,
      "save_acf_images":    true,
      "reanalyze":          false
    },
    "fit_lpi": {
      "enabled":   true,
      "channel":   "zenith-l",
      "max_dt":    300,
      "range_avg": [1, 3, 5],
      "reanalyze": false
    },
    "long_pulse": {
      "enabled": false,
      "channel": "zenith-l",
      "mode":    300
    },
    "fit_lp": {
      "enabled": false,
      "channel": "zenith-l",
      "avg_dur": 600,
      "ridx":    [35, 230]
    }
  }
}
\end{verbatim}

\subsubsection*{Key parameters to adjust}

\begin{description}
  \item[\texttt{channel}]
    Use \texttt{"zenith-l"} for the zenith antenna or \texttt{"misa-l"}
    for the steerable MISA antenna.
    Check which channels are present in \texttt{rf\_data/} before running.
  \item[\texttt{range\_gate\_us}]
    Range resolution in microseconds; must match the radar pulse-bit length.
    Typical values: 30, 60, or 120~\textmu s.
  \item[\texttt{max\_dt}]
    Integration time for the plasma-parameter fit.
    300~s is a good starting point; increase to 600~s for lower-SNR data.
  \item[\texttt{max\_time\_s}]
    Set to e.g.\ 600 for a quick 10-minute test run; remove the key
    (or set to \texttt{null}) to process the full experiment.
  \item[\texttt{reanalyze}]
    Set to \texttt{false} to skip intervals already on disk.
    Set to \texttt{true} to reprocess from scratch.
\end{description}

%-------------------------------------------------------------------
\subsection{Output Files}

All results are written to \texttt{output\_dir} (or inside the raw-data
directory if \texttt{output\_dir} is not set):

\begin{verbatim}
lpi_<range_gate_us>/zenith-l/
  lpi-<unix_timestamp>.h5    # ACF per range gate and lag (step 1)
  lpi-<unix_timestamp>.png   # diagnostic image

lpi_<range_gate_us>/zenith-l/
  pp-<unix_timestamp>.h5     # Te, Ti, vi, ne profiles (step 2)
  pp-<unix_timestamp>.png    # diagnostic plot

range_doppler_<mode>_outlier/zenith-l/
  il-<unix_timestamp>.h5     # range-Doppler spectra (step 3)
  pp-<unix_timestamp>.h5     # Te, Ti, vi, ne profiles (step 4)
\end{verbatim}

The plasma-parameter HDF5 files contain the following fields:
\begin{center}
\begin{tabular}{ll}
\hline
\textbf{Dataset} & \textbf{Content} \\
\hline
\texttt{Te}, \texttt{Ti}     & Electron / ion temperature (K) \\
\texttt{vi}                  & Line-of-sight ion drift velocity (m/s) \\
\texttt{ne}                  & Relative electron density proxy \\
\texttt{dTe/Ti}, \texttt{dTi}, \texttt{dvi}, \texttt{dne} & $1\sigma$ uncertainties \\
\texttt{rgs}                 & Range gate centres (km) \\
\texttt{t0}, \texttt{t1}     & Integration window (Unix seconds) \\
\texttt{az}, \texttt{el}     & Mean antenna pointing (degrees) \\
\texttt{T\_sys}              & System temperature (K) \\
\texttt{P\_tx}               & Mean transmit power (W) \\
\hline
\end{tabular}
\end{center}
% end of appendix


\section{Accessing the Analysis Server and Running the ISR Pipeline}
\label{app:analysis_guide}

This appendix describes how to obtain an account on the analysis server,
access a remote desktop session, retrieve the analysis code, set up the
Python environment, and run the Millstone Hill ISR plasma-parameter pipeline.

%-------------------------------------------------------------------
\subsection{Getting an Account on Revontuli}

The analysis is performed on \texttt{revontuli.uit.no}, a Linux server at
UiT -- The Arctic University of Norway.
Contact the group supervisor to have an account created.
You will receive a username and an initial password.
Change the password after first login:
\begin{verbatim}
ssh <username>@revontuli.uit.no
passwd
\end{verbatim}

%-------------------------------------------------------------------
\subsection{Remote Desktop Access}

A graphical desktop (XFCE4) is available via VNC.
Each user's VNC port is determined by their Unix user ID:
\begin{equation}
  \text{port} = \text{UID} + 50000
\end{equation}
Ask the system administrator for your UID, or read it yourself after
logging in with \texttt{id -u}.

\subsubsection*{Starting the desktop}

Log in over SSH and run the \texttt{remotedesktop} command (defined
system-wide in \texttt{/etc/profile.d/remotedesktop.sh}):
\begin{verbatim}
ssh <username>@revontuli.uit.no
remotedesktop
\end{verbatim}
This starts an Xtightvnc session with XFCE4 on your personal port.
To reconnect after a reboot, simply run \texttt{remotedesktop} again.

\subsubsection*{Connecting from your laptop}

Forward the VNC port over SSH so the connection is encrypted:
\begin{verbatim}
ssh -L 5901:localhost:<your-vnc-port> <username>@revontuli.uit.no
\end{verbatim}
Then open any VNC client (e.g.\ Remmina, TigerVNC, RealVNC) and connect
to \texttt{localhost:5901}.
The desktop appears without requiring a VNC password.

%-------------------------------------------------------------------
\subsection{Getting the Analysis Code}

The pipeline lives in a Git repository.
Clone it into your home directory on Revontuli:
\begin{verbatim}
cd ~
git clone https://github.com/jvierine/isr_analysis.git src/antistarlink
cd src/antistarlink
\end{verbatim}
To update to the latest version at any time:
\begin{verbatim}
git pull
\end{verbatim}

%-------------------------------------------------------------------
\subsection{Setting up the Python Environment}

A setup script installs all required Python packages into a virtual
environment at \texttt{\textasciitilde/venv/antistarlink}.
Run it once per machine:
\begin{verbatim}
cd ~/src/antistarlink
bash setup_env.sh
\end{verbatim}
The script requires \texttt{libfftw3-dev} to build \texttt{pyfftw}
(already installed on Revontuli).
It installs:
\begin{itemize}
  \item \texttt{pyfftw} -- fast FFT via FFTW3
  \item \texttt{digital\_rf} -- reading raw DigitalRF voltage data
  \item \texttt{mpi4py} -- MPI parallelisation
\end{itemize}
and inherits \texttt{h5py}, \texttt{scipy}, \texttt{numpy}, and
\texttt{matplotlib} from the system Python~3.14 installation.

%-------------------------------------------------------------------
\subsection{Raw Data Layout}

Each experiment directory contains the raw DigitalRF data and radar metadata:
\begin{verbatim}
/mnt/data/juha/millstone_hill/isr/<date>/
  usrp-rx0-r_<start>_<end>/
    rf_data/
      zenith-l/        # receive channel (1 MHz, complex64)
      misa-l/          # MISA receive channel (when available)
      tx-h/            # transmitted waveform copy
    metadata/
      id_metadata/     # pulse sweep IDs (coded=1-32, LP=300, scan=800)
      powermeter/      # transmit power vs. time
      antenna_control_metadata/   # pointing (azimuth, elevation)
\end{verbatim}

%-------------------------------------------------------------------
\subsection{Running the Pipeline}

The driver script \texttt{run\_analysis.py} reads a JSON configuration file
and executes the enabled pipeline steps in order.
Launch it with MPI to process intervals in parallel:
\begin{verbatim}
cd ~/src/antistarlink
mpirun -np 24 ~/venv/antistarlink/bin/python3 \
    run_analysis.py config/millstone_2023-09-05.json
\end{verbatim}
Replace \texttt{24} with the number of CPU cores available
(\texttt{nproc} reports the total on the current machine).

\subsubsection*{Configuration file}

An annotated example for the 2023-09-05 dataset is given below.
Create a copy and edit it for a new experiment.

\begin{verbatim}
{
  "experiment": "millstone_2023-09-05",
  "data_dir":   "/mnt/data/juha/millstone_hill/isr/2023-09-05/
                  usrp-rx0-r_20230905T214448_20230906T040054",
  "output_dir": "/mnt/data/juha/millstone_hill/results/2023-09-05",
  "max_time_s": 600,

  "steps": {
    "lpi": {
      "enabled":            true,
      "channel":            "zenith-l",
      "range_gate_us":      60,
      "avg_dur":            10,
      "pass_band_hz":       100000,
      "filter_len":         20,
      "max_range_delay_us": 7000,
      "save_acf_images":    true,
      "reanalyze":          false
    },
    "fit_lpi": {
      "enabled":   true,
      "channel":   "zenith-l",
      "max_dt":    300,
      "range_avg": [1, 3, 5],
      "reanalyze": false
    },
    "long_pulse": {
      "enabled": false,
      "channel": "zenith-l",
      "mode":    300
    },
    "fit_lp": {
      "enabled": false,
      "channel": "zenith-l",
      "avg_dur": 600,
      "ridx":    [35, 230]
    }
  }
}
\end{verbatim}

\subsubsection*{Key parameters to adjust}

\begin{description}
  \item[\texttt{channel}]
    Use \texttt{"zenith-l"} for the zenith antenna or \texttt{"misa-l"}
    for the steerable MISA antenna.
    Check which channels are present in \texttt{rf\_data/} before running.
  \item[\texttt{range\_gate\_us}]
    Range resolution in microseconds; must match the radar pulse-bit length.
    Typical values: 30, 60, or 120~\textmu s.
  \item[\texttt{max\_dt}]
    Integration time for the plasma-parameter fit.
    300~s is a good starting point; increase to 600~s for lower-SNR data.
  \item[\texttt{max\_time\_s}]
    Set to e.g.\ 600 for a quick 10-minute test run; remove the key
    (or set to \texttt{null}) to process the full experiment.
  \item[\texttt{reanalyze}]
    Set to \texttt{false} to skip intervals already on disk.
    Set to \texttt{true} to reprocess from scratch.
\end{description}

%-------------------------------------------------------------------
\subsection{Output Files}

All results are written to \texttt{output\_dir} (or inside the raw-data
directory if \texttt{output\_dir} is not set):

\begin{verbatim}
lpi_<range_gate_us>/zenith-l/
  lpi-<unix_timestamp>.h5    # ACF per range gate and lag (step 1)
  lpi-<unix_timestamp>.png   # diagnostic image

lpi_<range_gate_us>/zenith-l/
  pp-<unix_timestamp>.h5     # Te, Ti, vi, ne profiles (step 2)
  pp-<unix_timestamp>.png    # diagnostic plot

range_doppler_<mode>_outlier/zenith-l/
  il-<unix_timestamp>.h5     # range-Doppler spectra (step 3)
  pp-<unix_timestamp>.h5     # Te, Ti, vi, ne profiles (step 4)
\end{verbatim}

The plasma-parameter HDF5 files contain the following fields:
\begin{center}
\begin{tabular}{ll}
\hline
\textbf{Dataset} & \textbf{Content} \\
\hline
\texttt{Te}, \texttt{Ti}     & Electron / ion temperature (K) \\
\texttt{vi}                  & Line-of-sight ion drift velocity (m/s) \\
\texttt{ne}                  & Relative electron density proxy \\
\texttt{dTe/Ti}, \texttt{dTi}, \texttt{dvi}, \texttt{dne} & $1\sigma$ uncertainties \\
\texttt{rgs}                 & Range gate centres (km) \\
\texttt{t0}, \texttt{t1}     & Integration window (Unix seconds) \\
\texttt{az}, \texttt{el}     & Mean antenna pointing (degrees) \\
\texttt{T\_sys}              & System temperature (K) \\
\texttt{P\_tx}               & Mean transmit power (W) \\
\hline
\end{tabular}
\end{center}
% end of appendix


\section{Accessing the Analysis Server and Running the ISR Pipeline}
\label{app:analysis_guide}

This appendix describes how to obtain an account on the analysis server,
access a remote desktop session, retrieve the analysis code, set up the
Python environment, and run the Millstone Hill ISR plasma-parameter pipeline.

%-------------------------------------------------------------------
\subsection{Getting an Account on Revontuli}

The analysis is performed on \texttt{revontuli.uit.no}, a Linux server at
UiT -- The Arctic University of Norway.
Contact the group supervisor to have an account created.
You will receive a username and an initial password.
Change the password after first login:
\begin{verbatim}
ssh <username>@revontuli.uit.no
passwd
\end{verbatim}

%-------------------------------------------------------------------
\subsection{Remote Desktop Access}

A graphical desktop (XFCE4) is available via VNC.
Each user's VNC port is determined by their Unix user ID:
\begin{equation}
  \text{port} = \text{UID} + 50000
\end{equation}
Ask the system administrator for your UID, or read it yourself after
logging in with \texttt{id -u}.

\subsubsection*{Starting the desktop}

Log in over SSH and run the \texttt{remotedesktop} command (defined
system-wide in \texttt{/etc/profile.d/remotedesktop.sh}):
\begin{verbatim}
ssh <username>@revontuli.uit.no
remotedesktop
\end{verbatim}
This starts an Xtightvnc session with XFCE4 on your personal port.
To reconnect after a reboot, simply run \texttt{remotedesktop} again.

\subsubsection*{Connecting from your laptop}

Forward the VNC port over SSH so the connection is encrypted:
\begin{verbatim}
ssh -L 5901:localhost:<your-vnc-port> <username>@revontuli.uit.no
\end{verbatim}
Then open any VNC client (e.g.\ Remmina, TigerVNC, RealVNC) and connect
to \texttt{localhost:5901}.
The desktop appears without requiring a VNC password.

%-------------------------------------------------------------------
\subsection{Getting the Analysis Code}

The pipeline lives in a Git repository.
Clone it into your home directory on Revontuli:
\begin{verbatim}
cd ~
git clone https://github.com/jvierine/isr_analysis.git src/antistarlink
cd src/antistarlink
\end{verbatim}
To update to the latest version at any time:
\begin{verbatim}
git pull
\end{verbatim}

%-------------------------------------------------------------------
\subsection{Setting up the Python Environment}

A setup script installs all required Python packages into a virtual
environment at \texttt{\textasciitilde/venv/antistarlink}.
Run it once per machine:
\begin{verbatim}
cd ~/src/antistarlink
bash setup_env.sh
\end{verbatim}
The script requires \texttt{libfftw3-dev} to build \texttt{pyfftw}
(already installed on Revontuli).
It installs:
\begin{itemize}
  \item \texttt{pyfftw} -- fast FFT via FFTW3
  \item \texttt{digital\_rf} -- reading raw DigitalRF voltage data
  \item \texttt{mpi4py} -- MPI parallelisation
\end{itemize}
and inherits \texttt{h5py}, \texttt{scipy}, \texttt{numpy}, and
\texttt{matplotlib} from the system Python~3.14 installation.

%-------------------------------------------------------------------
\subsection{Raw Data Layout}

Each experiment directory contains the raw DigitalRF data and radar metadata:
\begin{verbatim}
/mnt/data/juha/millstone_hill/isr/<date>/
  usrp-rx0-r_<start>_<end>/
    rf_data/
      zenith-l/        # receive channel (1 MHz, complex64)
      misa-l/          # MISA receive channel (when available)
      tx-h/            # transmitted waveform copy
    metadata/
      id_metadata/     # pulse sweep IDs (coded=1-32, LP=300, scan=800)
      powermeter/      # transmit power vs. time
      antenna_control_metadata/   # pointing (azimuth, elevation)
\end{verbatim}

%-------------------------------------------------------------------
\subsection{Running the Pipeline}

The driver script \texttt{run\_analysis.py} reads a JSON configuration file
and executes the enabled pipeline steps in order.
Launch it with MPI to process intervals in parallel:
\begin{verbatim}
cd ~/src/antistarlink
mpirun -np 24 ~/venv/antistarlink/bin/python3 \
    run_analysis.py config/millstone_2023-09-05.json
\end{verbatim}
Replace \texttt{24} with the number of CPU cores available
(\texttt{nproc} reports the total on the current machine).

\subsubsection*{Configuration file}

An annotated example for the 2023-09-05 dataset is given below.
Create a copy and edit it for a new experiment.

\begin{verbatim}
{
  "experiment": "millstone_2023-09-05",
  "data_dir":   "/mnt/data/juha/millstone_hill/isr/2023-09-05/
                  usrp-rx0-r_20230905T214448_20230906T040054",
  "output_dir": "/mnt/data/juha/millstone_hill/results/2023-09-05",
  "max_time_s": 600,

  "steps": {
    "lpi": {
      "enabled":            true,
      "channel":            "zenith-l",
      "range_gate_us":      60,
      "avg_dur":            10,
      "pass_band_hz":       100000,
      "filter_len":         20,
      "max_range_delay_us": 7000,
      "save_acf_images":    true,
      "reanalyze":          false
    },
    "fit_lpi": {
      "enabled":   true,
      "channel":   "zenith-l",
      "max_dt":    300,
      "range_avg": [1, 3, 5],
      "reanalyze": false
    },
    "long_pulse": {
      "enabled": false,
      "channel": "zenith-l",
      "mode":    300
    },
    "fit_lp": {
      "enabled": false,
      "channel": "zenith-l",
      "avg_dur": 600,
      "ridx":    [35, 230]
    }
  }
}
\end{verbatim}

\subsubsection*{Key parameters to adjust}

\begin{description}
  \item[\texttt{channel}]
    Use \texttt{"zenith-l"} for the zenith antenna or \texttt{"misa-l"}
    for the steerable MISA antenna.
    Check which channels are present in \texttt{rf\_data/} before running.
  \item[\texttt{range\_gate\_us}]
    Range resolution in microseconds; must match the radar pulse-bit length.
    Typical values: 30, 60, or 120~\textmu s.
  \item[\texttt{max\_dt}]
    Integration time for the plasma-parameter fit.
    300~s is a good starting point; increase to 600~s for lower-SNR data.
  \item[\texttt{max\_time\_s}]
    Set to e.g.\ 600 for a quick 10-minute test run; remove the key
    (or set to \texttt{null}) to process the full experiment.
  \item[\texttt{reanalyze}]
    Set to \texttt{false} to skip intervals already on disk.
    Set to \texttt{true} to reprocess from scratch.
\end{description}

%-------------------------------------------------------------------
\subsection{Output Files}

All results are written to \texttt{output\_dir} (or inside the raw-data
directory if \texttt{output\_dir} is not set):

\begin{verbatim}
lpi_<range_gate_us>/zenith-l/
  lpi-<unix_timestamp>.h5    # ACF per range gate and lag (step 1)
  lpi-<unix_timestamp>.png   # diagnostic image

lpi_<range_gate_us>/zenith-l/
  pp-<unix_timestamp>.h5     # Te, Ti, vi, ne profiles (step 2)
  pp-<unix_timestamp>.png    # diagnostic plot

range_doppler_<mode>_outlier/zenith-l/
  il-<unix_timestamp>.h5     # range-Doppler spectra (step 3)
  pp-<unix_timestamp>.h5     # Te, Ti, vi, ne profiles (step 4)
\end{verbatim}

The plasma-parameter HDF5 files contain the following fields:
\begin{center}
\begin{tabular}{ll}
\hline
\textbf{Dataset} & \textbf{Content} \\
\hline
\texttt{Te}, \texttt{Ti}     & Electron / ion temperature (K) \\
\texttt{vi}                  & Line-of-sight ion drift velocity (m/s) \\
\texttt{ne}                  & Relative electron density proxy \\
\texttt{dTe/Ti}, \texttt{dTi}, \texttt{dvi}, \texttt{dne} & $1\sigma$ uncertainties \\
\texttt{rgs}                 & Range gate centres (km) \\
\texttt{t0}, \texttt{t1}     & Integration window (Unix seconds) \\
\texttt{az}, \texttt{el}     & Mean antenna pointing (degrees) \\
\texttt{T\_sys}              & System temperature (K) \\
\texttt{P\_tx}               & Mean transmit power (W) \\
\hline
\end{tabular}
\end{center}
% end of appendix
